\section{Evaluation}
\label{sec:evaluation}

The evaluation separates three objects that have different scopes.  The
release certificate binds the proof, statement, verifier binary, and concrete
security calculations.  Local validator runs measure both supported account
paths and exercise adversarial cases.  A mainnet-beta run records one
finalized verification-and-state-transition transaction for that exact
release.  Deployment and application setup precede the measured transition;
resource recovery and the treasury sweep follow it.

\subsection{Release instance}
\label{sec:evaluation-release}

The fail-closed release evaluator rebuilt the manifest-default SBF binary,
verified the concrete proof with the production host verifier, reconstructed
all three selector branches, reconciled the local compute ledger, and checked
36 required gates.  All gates passed and the failed-gate set was empty.  The
release-authorizing soundness value is the conservative whole-ledger,
three-branch floor of 100.16144938287457 bits.  The sharper selector-scoped
calculation, 101.38010539241552 bits, is retained as a diagnostic and is not
used for the headline claim.

\begin{table}[t]
  \centering
  \small
  \caption{Concrete \Profile{} release instance.}
  \label{tab:release-schedule}
  \begin{tabular}{l l}
    \toprule
    Item & Value \\
    \midrule
    trace rows / code rate & $2^{10}$ / $1/512$ \\
    query fibers & $q=18$ \\
    batch work & 37 bits \\
    fold work & $(34,33,30,25)$ bits \\
    final work & 32 bits \\
    selector candidates / selected & 3 / 0 \\
    proving-attempt cap & 17 \\
    release gates & 36 / 36 \\
    conservative soundness floor & 100.161449 bits \\
    real-versus-simulator floor & 104.024922 bits \\
    pairwise-witness floor & 103.024922 bits \\
    \bottomrule
  \end{tabular}
\end{table}

The content identities used throughout the evaluation are:
\begin{description}
  \item[Proof.] \ProfileProofBytes{} bytes; SHA-256
  \nolinkurl{d4f529964d1cf9ccd9c5568b694796ba54191c6be38d341c66efa08c830cdc3d}.
  \item[Statement sidecar.] 667 bytes; SHA-256
  \nolinkurl{947a608c93487a634f37119bead8d61fe29e9cb6883493465d6fb35af27883c2}.
  \item[Canonical public-input digest.] SHA-256
  \nolinkurl{b2d150dfcb6432c1b6f2e3892ee45a9aa5f393809d97c8292fea975b3da35fa3}.
  \item[Production SBF.] \ProfileSbfBytes{} bytes; SHA-256
  \nolinkurl{97c45a9abef97607a2fc6ed245829210046b234044b6738599d2bce0c367d04a}.
  \item[Release certificate.] SHA-256
  \nolinkurl{4ace42c0954bee74382520bbb0818da86e570ec9b212c61cb21130dc644a2725}.
\end{description}

The statement has sequence zero, asset identifier 17, and fee 1.  Each of the
three reconstructed schedules met the required $\mathsf{Good}_{23}$ ranks,
so selector zero was both serialized and the least accepted selector.  The
proof, statement, and SBF identities were required to agree across host
acceptance, local mutation, release, and mainnet preflight artifacts.
The execution-time preflight used certificate
\nolinkurl{d83f6df645a00f28090e78f8648974ec94afd35f4677f5d84525ba6af7a2688f}.
The publication certificate was regenerated after correcting stale,
non-authorizing release metadata in its soundness source, removing duplicated
release snapshots from the proof-independent complete-Good and HVZK sources,
replacing local absolute command paths in the acceptance and mutation sources
with repository-relative paths, and correcting a q18 comment in the program
manifest.  The explanatory certificate note was updated to describe that
non-authorizing metadata boundary.  A machine diff is limited to that note,
the generation timestamp, byte lengths and digests for the five corrected
JSON sources, and the program-manifest digest.  The proof, statement, SBF,
security values, compute values, and all 36 gate outcomes are identical.  The
release bundle retains both certificates.

\subsection{Local compute and functional tests}
\label{sec:evaluation-local}

Local compute was measured by transaction simulation on
\texttt{solana-test-validator 2.3.0} (Agave source identifier
\texttt{a2e21dda}, feature set \texttt{3640012085}).  The frozen SBF was built
with \nolinkurl{solana-cargo-build-sbf} version \texttt{2.3.0}, platform-tools
\texttt{v1.48}, and its bundled Rust \texttt{1.84.1}; a build producing
different bytes fails the release identity gate.  Each transaction
requested a 1,400,000-CU limit and a 262,144-byte heap.  The read-only tag~59
baseline and both tag~65 paths used the same uninstrumented release binary.
Table~\ref{tab:production-cu} reports literal transaction totals.

\begin{table}[t]
  \centering
  \small
  \caption{Production-binary local compute consumption.}
  \label{tab:production-cu}
  \begin{tabular}{l r r r}
    \toprule
    Marker path & Tag~59 CU & Tag~65 CU & Headroom \\
    \midrule
    program-owned zero marker & 1,303,642 & 1,338,471 & 61,529 \\
    System-owned marker creation & 1,303,642 & 1,340,803 & 59,197 \\
    \bottomrule
  \end{tabular}
\end{table}

The local production matrix also admitted only the intended feature set and
rejected combinations with diagnostic and superseded profiles.  It accepted
the mined proof on tags 59 and 65, rejected a separately committed unmined
proof, and submitted corrupted-proof transactions to check rollback.  Both
marker paths produced the expected pool and nullifier images; a duplicate was
rejected without a second mutation.  A two-signer race on the canonical
System-owned marker path produced exactly one commit.  Successful tag~65 runs
closed the proof account and reconciled its complete lamport refund.

\subsection{Mainnet-beta method}
\label{sec:evaluation-mainnet-method}

The mainnet run used the Solana genesis hash
\begin{center}
  \small\ttfamily
  5eykt4UsFv8P8NJdTREpY1vzqKqZKvdpKuc147dw2N9d.
\end{center}
A read-only finalized snapshot first required the disposable Program,
ProgramData, pool, and nullifier addresses to be absent.  The executor then
deployed the exact release binary under program
\begin{center}
  \small\ttfamily
  9kPpUknrRicMvaGa6zPNERGUYDj6fMvMR8PwMS3iFR6Z
\end{center}
with loader-derived ProgramData address
\begin{center}
  \small\ttfamily
  G7XHkv6riJ8TQJ6R8o31wAUMaoL8Q5UyHGExpQ51ntp1.
\end{center}
The finalized deployment transaction was recorded at slot 432,933,454.  The
executor checked the deployed code hash, exact maximum data length, loader
ownership, Program--ProgramData linkage, and the dedicated upgrade authority
before application setup, immediately before verification simulation, and
after verification finality.

Application setup comprised \ProfileMainnetSetupTransactions{} transactions:
pool creation and initialization, combined proof-account creation and tag~0
initialization, \ProfileUploadTransactions{} 960-byte-or-smaller upload
transactions, and tag~62 finalization.  Uploads were submitted in five
finality windows of at most 16 transactions.  The complete 64,487-byte proof
account image was compared with the expected image before sealing.  Simulated
post-finalization upload and second-finalization attempts were rejected.
Program deployment and these application-setup transactions are not part of
the one-transaction verification claim.

The executor constructed one tag~65 transaction with a fresh blockhash,
simulated the exact signed bytes with signature verification and without
blockhash replacement, and submitted those same bytes.  It allowed no changed
transaction after signing and used no identical-wire retry.  Finalized
account images, transaction metadata, program continuity, and the program-log
proof digest were checked before the run was recorded as successful.

\subsection{Finalized mainnet-beta verification transition}
\label{sec:evaluation-mainnet-result}

The tag~65 transaction finalized at slot \ProfileMainnetSlot{} on 14 July
2026 at 21:40:36 UTC.  Its signature is
\begin{center}
  \scriptsize\ttfamily
  \href{https://explorer.solana.com/tx/4Er5afhxfcFmpeTuFqEeNQEbCBri3pkc6ymx7ST5wfNpSBYwQDHA9DtCuDBpD5WuEDXs7ozL3sK5msc6QWE4q9Fo?cluster=mainnet-beta}{4Er5afhxfcFmpeTuFqEeNQEbCBri3pkc6ymx7ST5wfNpSBYwQDHA9DtCuDBpD5WuEDXs7ozL3sK5msc6QWE4q9Fo}.
\end{center}
The signed wire had SHA-256
\nolinkurl{09183f6b5e95872b884cf549e7dc5fecd1862b0c7c5b77ed8eb7f57dcbf5f5db};
its message had SHA-256
\nolinkurl{7565f8cbe265d4be73c3f79d73b4bcc9c48d59a13f1ef15a031054a0f8836670}.
Simulation and the finalized record both reported
\ProfileMainnetComputeUnits{} CU.  This is 95.9821\% of the requested limit
and leaves \ProfileMainnetComputeHeadroom{} CU.  The transaction fee was
10,000 lamports and the metadata reported no error.

The program emitted the domain-separated proof-digest log
\texttt{aspis-proof-sha256-v1}.  Decoding its value gives exactly the release
proof SHA-256 listed in Section~\ref{sec:evaluation-release}; the executor
required this equality in both simulation and finalized logs.  The pool at
\nolinkurl{B1z4gQ82xerghDKZ68HuP2Tehx5ER5mfVHq74y4ew3kk} advanced from sequence
zero to one and its data hash changed from
\nolinkurl{3cad0b60e8a697c28f7c16657ed19142893f2eb5513dc0bc1382942d1789d11e}
to
\nolinkurl{3d259b401a8f6c5cb49713707a7108b4097e35873e893cea1a93bc46a5d78ea4}.
The previously absent canonical nullifier marker
\nolinkurl{2CnUJycxkinN3XtpWH5bCYVDEuMx42BTnFU9WXsymAu3} was created with the
expected 72-byte image.

The proof account held \ProfileProofRentLamports{} lamports immediately before
the transition and zero afterward.  The payer increased by the full proof
balance, less the 1,392,000 lamports used to fund the new nullifier and the
10,000-lamport transaction fee.  The evidence reconciles this equation from
the finalized pre- and post-balances.  Proof verification, proof-account
refund, nullifier creation, and pool mutation therefore occurred in the same
transaction; none of the preceding setup transactions is included in that
statement.

After finality, a duplicate-spend simulation used a separately created and
sealed replay-probe account.  It returned the exact nullifier-already-spent
error and left the pool, nullifier, and replay-probe images unchanged.  The
probe was then closed with tag~64 and its 1,176,240-lamport balance was
refunded.  These replay-probe transactions are post-transition tests, not part
of the verification transaction.

\subsection{Resource recovery and cost accounting}
\label{sec:evaluation-cleanup}

The deployment was deliberately disposable.  After the execution evidence
had been made read-only, a separate loader-v3 close transaction finalized at
slot 432,934,227.  It removed ProgramData and refunded
\ProfileProgramDataRefundLamports{} lamports; the exact refund equation and
10,000-lamport fee were reconciled from the landed transaction.  The small
loader Program account, pool, and nullifier marker remained.  Closing
ProgramData disabled the disposable deployment after the measured
verification; it did not alter the already finalized transaction record.

The later transfer of the remaining payer balance back to the funding source
is treasury cleanup and is not a protocol result.  Its separate finalized
receipt records 6,985,137,600 lamports returned at slot 432,934,316 and a zero
post-sweep payer balance.  Of the demo payer's balance, 0.0148834 SOL was not
returned: 0.0037584 SOL remained as rent in the Program, pool, and nullifier
accounts, and 0.0111250 SOL was paid as aggregate payer transaction fees.
These figures include deployment and every payer-funded setup and cleanup
transaction; they exclude the inbound funding transaction fee paid by the
funding source and are not the fee of the single verification transition.

\subsection{Independent reconciliation and reproducibility}
\label{sec:evaluation-reproduction}

The source snapshot, frozen publication bundle, and paper are available under
the immutable tag
\href{https://github.com/DJBarker87/zk/tree/aspis-spend-q18-g37-mainnet-v1}{\texttt{spend-q18-g37-mainnet-v1}}.
From that checkout, the bundle's offline verifier is
\path{release/aspis-spend-q18-g37-mainnet-v1/verify.sh}; it authenticates every
bundled object and checks the certificate lineage, proof and statement
identities, deployed ProgramData image, finalized transaction observations,
and cost equations.

The public reconciliation file
\path{results/spend/spend_mainnet_independent_rpc_reconciliation.json}
has SHA-256
\nolinkurl{72dcbc71d2657cd58a1af8d6f28e355949f63f67c4746e8d80573ecba159cabc}.
It records sanitized raw \texttt{getTransaction} responses from the official
Solana mainnet endpoint and an independent PublicNode endpoint.  Both report
the same mainnet genesis, signature, slot, block time, error status, fee,
compute consumption, balances, proof-digest log, and success log for the
verification transaction.  Both also agree on the funding, deployment,
replay-probe close, ProgramData close, and treasury-sweep records and on the
final account snapshot.  Queries for the verification signature on official
devnet and testnet returned null, which disambiguates the cluster.  This
machine-readable reconciliation, rather than an explorer rendering, is the
source for the network claims in this section.

The archival reproducer
\path{tools/reconstruct_spend_mainnet_sbf.py} independently obtains the
same finalized 1,069-signature buffer history and raw transaction bodies from
both providers.  It binds each of 1,067 loader writes to the pinned buffer and
authority and reconstructs all 921,848 SBF bytes without gaps or overlaps.
The recovered SHA-256 is
\nolinkurl{97c45a9abef97607a2fc6ed245829210046b234044b6738599d2bce0c367d04a},
equal to the bundled SBF.  Applying the loader-v3 ProgramData header gives
\nolinkurl{7fe4c7352f19ed00a3ba72b3dc6038943db3e5311c23cd5c1b0f8dbb15c4cda2},
equal to the pre-close account-data digest.  The same replay checks the
deployment identities and decodes the landed 137-byte tag-65 payload, account
order, proof-digest log, and refund equation.  Its frozen output is
\path{results/spend/spend_mainnet_sbf_and_instruction_reconstruction.json}.

The public finalized manifest is
\path{results/spend/spend_mainnet_finalized_manifest.json}, with SHA-256
\nolinkurl{a8501bee992c48d1777f01be78d8406580a5955e819d8c76ac4426ce7c03a604}.
It contains no local absolute or operational filesystem paths or key material
and binds the release certificate,
the immutable execution and cleanup receipts, the independent reconciliation,
all principal signatures, the rent refunds, and the final cost equations.

The release certificate can be reevaluated from the repository root with:
\begin{verbatim}
NO_DNA=1 cargo run -q -p aspis-prover --example \
  spend_soundness_epro_ledger -- --calculation-only
NO_DNA=1 cargo run --release -p aspis-xtask -- \
  spend-release
\end{verbatim}
The release command rebuilds the manifest-default SBF before comparing it
with the production-only binary and evaluating all 36 gates.  Reproduction
requires the proof and statement bytes; their hashes are identifiers, not
substitutes.  The evaluator records its current generation time, so a later
reevaluation reproduces the substantive fields and gate outcomes but not the
frozen certificate's byte identity or SHA-256.
